\chapter{РЕАЛИЗАЦИЯ ПРОГРАММНОГО КОМПЛЕКСА}
\label{ch:implementation}
\sloppy


\section{Архитектура реализованного программного комплекса}
\label{sec:3-1}

Глава~\ref{ch:method} определяет состав системы как набор из девяти функциональных модулей, образующих систему поддержки принятия решений для организатора программы конференции. Настоящий подраздел описывает реализацию этого замысла в виде программного комплекса.

Комплекс собран из нескольких модулей: модели, симуляторы, рекомендательные политики, оператор локальных модификаций программы, планировщик экспериментов, модуль показателей и средства формирования сравнительных отчётов. Модули взаимодействуют через стабильные интерфейсы; это даёт воспроизводимость отдельных этапов эксперимента, независимое тестирование модулей и расширение состава политик и моделей популярности без переработки симулятора.

В реализации комплекс разделён на три архитектурных уровня. Уровень данных хранит и подготавливает входы: программа конференции, синтетическая популяция профилей участников, векторные представления описаний докладов и профилей, преднасчитанная оценка априорной популярности докладов. Уровень моделирования реализует отклик аудитории на программу и на рекомендации: модель поведения участника, семейство политик рекомендаций, оператор локальных модификаций программы и два симулятора с различной формой модели выбора --- параметрический и агентский на основе языковой модели. Уровень эксперимента отвечает за дизайн численного эксперимента и обработку его результатов: генератор плана эксперимента, управление случайностью на принципе общих случайных чисел, модуль показателей загрузки и риска, средства формирования сравнительного отчёта. Зависимости между уровнями однонаправленны: данные не знают ничего о моделях и эксперименте; моделирование зависит от данных, но не от плана эксперимента; эксперимент собирает результаты, но не вмешивается в логику отклика. Программа конференции остаётся единственным фактическим входом системы; все остальные сущности либо моделируются, либо задаются параметрами эксперимента.

Разработанный программный комплекс рассматривается как аналитический слой существующей продуктовой экосистемы для конференций сообщества JUG, на стороне организатора. Клиентская часть экосистемы --- приложение для индивидуального формирования участником расписания докладов --- была развёрнута в производственной среде и обеспечивает существующий канал взаимодействия с участниками. Комплекс подключается к этой экосистеме как самостоятельный аналитический инструмент: на вход он принимает программу конференции в формате, совместимом с существующим сервисом рекомендаций, и возвращает сравнительный отчёт для организатора. Метрики использования клиентского приложения и пользовательские сигналы в постановку настоящей работы не входят и комплексом не используются.

Реализация программного комплекса выполнена автором лично. К авторским инженерным решениям относятся: разделение каналов влияния политики и состояния системы на функцию полезности участника; форма функции полезности с симплексным соотношением весов и реализованные представители её компонентов; форма оператора локальных модификаций программы и протокол его применения; паритет каналов между параметрическим и агентским симуляторами как принцип проектирования контекста для языковой модели; схема разделения источников случайности при общих случайных числах; протокол прохождения граничной верификации модели как блокирующего фильтра перед содержательными выводами. Существующий клиентский сегмент продуктовой экосистемы JUG используется как контекст интеграции и в состав разработанного комплекса не входит.

Три архитектурных решения определяют комплекс в целом и обсуждаются ниже. Первое --- разделение каналов влияния политики и состояния системы на функцию полезности участника: политика влияет на распределение посещений строго через компонент функции полезности с весом $w_{rec}$, а ограничения вместимости учитываются на уровне политики, не на уровне выбора участника. Благодаря этому требование о граничном поведении модели --- при $w_{rec} \to 0$ политики неразличимы --- выполняется по построению, а риск двойного учёта вместимости снят. Подробное обсуждение реализации модели поведения вынесено в подраздел~\ref{sec:3-3}.

Второе решение --- двусимуляторная архитектура моделирования отклика. Параметрический симулятор работает в замкнутой форме выбора и используется как основной инструмент полного прогона по плану эксперимента; агентский симулятор на основе языковой модели работает на репрезентативном подмножестве конфигураций и используется для перекрёстной проверки выводов о сравнительной устойчивости политик. Оба симулятора реализованы как взаимозаменяемые компоненты, работающие через единый интерфейс с модулем программы и модулем политик. Это удовлетворяет защищаемому положению о взаимной проверке моделей; реализация описана в подразделе~\ref{sec:3-5}.

Третье решение --- построение плана эксперимента на основе латинского гиперкуба с применением общих случайных чисел внутри одной точки выборки. План эксперимента отделён от логики симуляции и реализован как самостоятельный модуль, формирующий выборку конфигураций по заданным параметрическим осям. При общих случайных числах внутри одной точки гиперкуба все политики применяются к одной и той же синтетической аудитории и одной и той же реализации варианта программы; шум сравнения политик при этом существенно снижается. Реализация описана в подразделе~\ref{sec:3-6}.

Дальнейшие подразделы настоящей главы последовательно раскрывают сценарий выполнения комплекса (подраздел~\ref{sec:3-2}), реализацию модели отклика участника (подраздел~\ref{sec:3-3}), реализацию управляющих воздействий --- политик рекомендаций и оператора локальных модификаций программы (подраздел~\ref{sec:3-4}), агентский симулятор как второй источник отклика (подраздел~\ref{sec:3-5}), реализацию плана эксперимента и обработки результатов (подраздел~\ref{sec:3-6}), использованные данные и верификацию работоспособности (подраздел~\ref{sec:3-7}), технологический стек и порядок использования систем искусственного интеллекта (подраздел~\ref{sec:3-8}). Подраздел~\ref{sec:3-9} содержит выводы по главе.


\section{Пайплайн выполнения сценарного эксперимента}
\label{sec:3-2}

Сценарный эксперимент устроен как набор однотипных единичных итераций. Единичная итерация --- вычисление показателей загрузки, риска перегрузки и релевантности для одной комбинации <<программа конференции --- конфигурация модели --- политика рекомендаций --- реализация случайности>>. Подготовка данных, перебор конфигураций и агрегация результатов обслуживают этот единичный прогон.

На входе фиксируется программа конференции --- структурированное описание слотов, докладов, залов и их номинальных вместимостей, а также метаданные докладов: тексты аннотаций, тематические признаки, состав спикеров. Номинальная вместимость отражает физические свойства площадки и сама по себе в исходных данных не варьируется; её масштабирование в сценарных конфигурациях выполняется отдельной осью эксперимента, описанной ниже, и применяется в момент прогона без изменения исходной программы. Программа загружается из файла фиксированного формата; векторные представления описаний докладов и профилей участников считаются заранее, на этапе подготовки данных, и хранятся отдельно. Это разделение продиктовано двумя соображениями. Первое --- векторные представления получены однократным вызовом внешней модели представления текстов и не должны пересчитываться при каждом запуске эксперимента. Второе --- отделение представлений от программы делает возможным независимую замену источника эмбеддингов без изменения логики симулятора.

Параметры сценарного эксперимента формируются генератором плана эксперимента. План задаёт набор точек в пространстве параметрических осей конфигурации, перечисленных в главе~\ref{ch:method}: вместимость залов, источник оценки популярности докладов, веса модели поведения участника, размер аудитории, вариант программы. Внутри каждой точки гиперкуба перебираются все четыре политики рекомендаций; каждая комбинация <<точка $\times$ политика>> повторяется на нескольких независимых случайных зёрнах для усреднения. План построен на латинском гиперкубе, дающем равномерное покрытие пространства параметров заданным числом точек при существенном сокращении объёма по сравнению с полным факторным перебором. Подробное обсуждение этого решения вынесено в подраздел~\ref{sec:3-6}.

Перед запуском единичной итерации система формирует синтетическую аудиторию. Аудитория получается отбором из заранее подготовленной популяции профилей: это сокращает стоимость формирования и обеспечивает воспроизводимость. Размер аудитории задаётся осью эксперимента, состав отбора --- отдельным случайным зерном, фиксированным на уровне точки гиперкуба и одинаковым для всех политик внутри этой точки. Сравнение политик внутри одной точки выполняется на одной и той же синтетической аудитории. Это реализация принципа общих случайных чисел из методики планирования имитационных экспериментов; она снимает шум сравнения политик и делает попарные сопоставления внутри точки информативными.

Источник оценки популярности докладов выбирается параметром конфигурации. Реализованы три варианта: оценка только по содержательной близости описания доклада к тематике конференции, оценка только по структурным признакам известности доклада в программе и комбинированная оценка. Соответствие источника популярности и поведения участника задаётся через функцию полезности и фиксируется в подразделе~\ref{sec:3-3}.

Вариант программы --- это либо исходная программа, либо одна из её допустимых локальных модификаций, полученных оператором перестановок докладов между разными временными слотами одного дня. Реализация оператора и проверка ограничения о конфликтах спикеров описаны в подразделе~\ref{sec:3-4}. Множество допустимых вариантов программы фиксируется на уровне точки гиперкуба и одинаково для всех политик внутри этой точки; это продолжает принцип общих случайных чисел и делает сопоставимыми результаты сравнения политик при одной и той же реализации структуры программы.

Единичный прогон симулятора обрабатывает программу слот за слотом. Внутри одного временного слота участники обрабатываются последовательно: каждый следующий участник наблюдает обновлённое состояние загрузки залов, отражающее решения тех, кто был обработан раньше. Это даёт реалистичный эффект кумулятивного формирования распределения посещений. Шаг для одного участника состоит из трёх связанных действий: политика формирует упорядоченный список рекомендаций; модель поведения участника вычисляет распределение вероятностей выбора над докладами слота; симулятор разыгрывает выбор и обновляет состояние загрузки. Логика модели поведения и логика политики разделены явно: политика влияет на выбор участника только через компонент функции полезности с соответствующим весом, как зафиксировано в постановке.

После прохождения всех слотов программы для одной комбинации параметров система фиксирует протокол шагов и состояние загрузки, накопленное к концу каждого слота. На основе этих данных модуль показателей вычисляет метрики из набора, заданного в главе~\ref{ch:method}. Показатели сохраняются в табличном формате с явным указанием значения каждой оси конфигурации и идентификатора политики, что даёт возможность последующего попарного сравнения и анализа чувствительности.

Воспроизводимость прогонов держится на строгом разделении источников случайности. Случайность в системе возникает в трёх местах: при отборе синтетической аудитории, при формировании множества вариантов программы оператором локальных модификаций и при разыгрывании выбора участника в симуляторе. Каждый из этих источников управляется отдельным случайным зерном с явным правилом образования. Зерно отбора аудитории и зерно оператора модификаций программы фиксируются по идентификатору точки гиперкуба и одинаковы для всех политик и реплик внутри этой точки; зерно симулятора фиксируется по номеру реплики и варьируется между ними, что и даёт независимое усреднение. Внутри самого симулятора разделены поток случайности окончательного выбора участника и поток случайности стохастической части политики. Это делает граничное поведение модели строго воспроизводимым: при нулевом весе компонента рекомендации выбор участника не зависит от политики --- свойство, которое глава~\ref{ch:method} относит к обязательным для граничной верификации.

Завершение единичной итерации даёт строку в табличном формате результатов. Полный сценарный эксперимент состоит из последовательного выполнения всех таких итераций по плану эксперимента и их объединения в единую таблицу. Дальнейшая обработка результатов --- попарные сравнения политик, оценка устойчивости при варьировании одной из осей конфигурации, эффект локальных модификаций программы относительно исходной программы и эффект социального заражения в зависимости от параметров --- выполняется отдельным модулем постобработки на уже сформированной таблице. Результаты этой обработки сохраняются в виде отдельных артефактов и подаются на вход модулю формирования сравнительного отчёта. Численные результаты постобработки и их интерпретация представлены в главе~\ref{ch:results}.


\section{Реализация модели отклика участника}
\label{sec:3-3}

Модель отклика участника --- центральный компонент параметрического симулятора. На каждом шаге симуляции она получает на вход профиль участника, множество доступных в текущем временном слоте докладов, упорядоченный список рекомендаций политики и текущее состояние загрузки залов; на выходе задаёт распределение вероятностей выбора над докладами слота, по которому симулятор разыгрывает фактический выбор. Все сравнительные результаты эксперимента определяются формой этой модели и значениями её параметров.

Модель реализована в форме мультиномиальной логит-модели с аддитивной функцией полезности из трёх компонентов --- содержательной релевантности, следования рекомендации и социального заражения; стохастичность выбора задаётся отдельным параметром температуры. Эти компоненты введены в постановке главы~\ref{ch:method}; здесь рассматривается их инженерная реализация. В принятой реализации в качестве канала следования рекомендации использована индикаторная форма, а в качестве компонента социального заражения --- логарифмическая нормировка по числу выбравших доклад участников; оба представителя выбраны как простейшие содержательные формы соответствующих классов и подробно обсуждаются ниже. Веса трёх компонентов связаны симплексным ограничением: их сумма равна единице. Симплексное ограничение реализовано на уровне генератора плана эксперимента, не на уровне модели поведения; модель поведения принимает три веса как независимые параметры. Это упрощает интерпретацию весов как долей влияния трёх каналов и сокращает размерность пространства параметрических осей, не теряя содержания постановки.

Реализация опирается на четыре инженерных решения, каждое из которых обсуждается ниже.

\paragraph{Учёт вместимости как ответственность политики, не модели поведения.}
В первой реализации параметрического симулятора, выполненной до фиксации текущей постановки, ограничение вместимости было зашито в функцию полезности участника как штрафной член, монотонно убывающий по текущей загрузке зала. Это решение пересмотрено по двум содержательным причинам.

Первая причина --- противоречие с принципом разделения каналов, принятым в главе~\ref{ch:method}: политика рекомендаций влияет на распределение посещений строго через компонент функции полезности с весом $w_{rec}$. Если ограничение вместимости входит в функцию полезности напрямую, оно действует в обход этого канала: даже контрольная политика, не предъявляющая участнику никакой рекомендации, оказывается косвенно подвержена эффекту вместимости. Это нарушает разделение каналов и делает невозможным выполнение свойства неразличимости политик при $w_{rec} \to 0$, обязательного по требованию граничной верификации модели.

Вторая причина --- расхождение с практикой, наблюдаемой в литературе. В рассмотренных в главе~\ref{ch:domain} работах по рекомендательным системам с учётом ограничений вместимости вместимость трактуется как свойство стороны рекомендатора, а не стороны участника. Метод ReCon реализует учёт через целевую функцию обучения рекомендатора с компонентом транспортной стоимости; формализм FEIR --- через переранжирование выдачи; capacity-aware методы в области рекомендации точек интереса --- через постпроцессинговый шаг распределения. Во всех трёх случаях участник остаётся моделью индивидуального выбора без явного знания о вместимости.

В принятой реализации ограничение вместимости полностью изъято из функции полезности участника. Учёт вместимости целиком ложится на одну из политик семейства --- политику с учётом загрузки залов, описанную в подразделе~\ref{sec:3-4}. Эта политика добавляет к скорингу кандидатов компонент, монотонно убывающий по текущей загрузке зала, и отбрасывает залы с загрузкой выше заданного жёсткого порога. Так достигается симметричное соответствие постановке: политика влияет на выбор участника только через компонент полезности с весом $w_{rec}$, а ограничение вместимости становится одним из возможных дизайнов политики и сравнивается с альтернативами на общих основаниях.

\paragraph{Канал рекомендации как индикатор присутствия в выдаче.}
Канал следования рекомендации в принятой реализации задан индикаторной функцией: компонент полезности равен единице, если доклад входит в упорядоченный список рекомендаций, выданный политикой, и нулю в противном случае. Это самая простая из возможных форм, выбрана сознательно.

Альтернативные формы --- ранговая, в которой вклад убывает с позицией доклада в списке, и нормированная по скорингу политики --- рассматривались на этапе проектирования модели. Содержательно ни одна из них не была отвергнута, но в условиях отсутствия данных о фактическом поведении аудитории выбор более сложной формы означает введение в модель дополнительной свободы без возможности её калибровки. Индикаторная форма зафиксирована как минимальная содержательная: она удовлетворяет требованию о граничном поведении модели по $w_{rec}$ и оставляет возможность последующего расширения без переписывания общей структуры функции полезности.

В контрольной политике, которая по постановке не предъявляет участнику никакой рекомендации, индикаторный компонент тождественно равен нулю для всех докладов слота. В этом случае функция полезности сводится к взвешенной сумме компонента релевантности и компонента социального заражения, и распределение выбора определяется содержательной близостью докладов к профилю участника и кумулятивным эффектом в когорте. В этом и состоит роль контрольной политики как точки отсчёта.

\paragraph{Социальное заражение как логарифмический сигнал по числу выбравших доклад.}
Компонент социального заражения отражает динамическое усиление выбора в зависимости от уже сделанных выборов в когорте: формально это монотонно возрастающая функция от числа участников, уже выбравших данный доклад к моменту принятия решения текущим участником. На этапе проектирования рассматривался ряд форм этой зависимости --- линейная пропорциональная по числу выбравших, насыщающая функция с экспоненциальным затуханием, нормированная доля внутри слота, нормировка относительно среднего по слоту, логарифмическая шкала, форма на основе распространения информационных каскадов и несколько производных от перечисленных.

В принятой реализации использована форма с логарифмической нормировкой числа выбравших к логарифму общего размера аудитории. Эта форма сочетает два полезных свойства: монотонный рост влияния с числом ранее сделанных выборов в когорте и насыщение при больших значениях, что предотвращает доминирование первого популярного доклада за счёт эффекта <<победитель забирает всё>>. Она не вводит в модель новых параметров, кроме веса самого компонента, и встречается в литературе по агентским симуляторам пользователей рекомендательных систем.

Значение компонента социального заражения вычисляется по числу выбравших данный доклад, а не по доле занятых мест в зале. Это архитектурно существенно. На текущих программах конференций, где в одном временном слоте каждый зал содержит ровно один доклад, две величины численно совпадают. Семантически они различны: число выбравших доклад --- это популярность доклада, доля занятых мест в зале --- ограничение вместимости. Первая величина даёт понятийное разделение каналов <<социальный сигнал>> и <<ограничение вместимости>> и поддерживает принятое решение о выносе ограничения вместимости в политику. При программах с другой структурой, где в одном зале в одном слоте может проходить несколько докладов, разделение сохраняется без переписывания модели.

В отличие от компонента следования рекомендации, компонент социального заражения вычисляется детерминированно по фиксированному состоянию счётчика выборов и не использует генератора случайных чисел. Это свойство существенно для воспроизводимости сравнения политик: при нулевом весе компонента социального заражения функция полезности сводится к двум остальным компонентам, и вся стохастическая логика прогона совпадает с реализацией без социального заражения с точностью до пословного протокола шагов. Тем самым требование совместимости с базовой моделью при выключенном социальном заражении выполняется по построению, не за счёт отдельной логики.

\paragraph{Распределение выбора и альтернатива отказа от посещения.}
Распределение вероятностей выбора над докладами слота получается применением функции softmax к полученным значениям полезности с делением на параметр температуры. Дополнительно в распределение явно введена альтернатива отказа от выбора с фиксированной вероятностью; остальная вероятностная масса распределяется между докладами слота. Альтернатива отказа моделирует наблюдаемую на конференциях долю участников, не посещающих текущий слот по причинам, не связанным с программой; явное введение <<no-choice alternative>> соответствует распространённой практике моделирования индивидуального выбора в литературе по дискретному выбору.

\paragraph{Простота модели и её сценарная роль.}
Принятая модель отклика проста по форме и интерпретируема по компонентам: три аддитивных канала, симплексное соотношение весов, общая логит-форма распределения. Эта простота продиктована содержательным ограничением задачи: данные о фактической посещаемости отсутствуют, и калибровка модели на эмпирических наблюдениях за индивидуальным поведением невозможна. Введение более тонкой структуры --- гетерогенности параметров по латентным классам, вложенной логит-формы с группированием близких альтернатив, моделей с динамической памятью на уровне отдельного участника --- означает добавление в модель свободы, которая не может быть подтверждена данными и превращается в источник трудно интерпретируемых артефактов. Допущение о параметрической форме модели выбора, на котором держится эта простота, зафиксировано в главе~\ref{ch:method}.

Принятая логит-форма удобна по нескольким причинам. Граничная верификация по весу компонента следования рекомендации выполняется по построению: при нулевом весе все политики становятся неразличимы по функции полезности. Режим общих случайных чисел сохраняется при попарном сравнении политик внутри точки гиперкуба. Состав компонентов функции полезности можно расширять, не переписывая её общей структуры; так в реализации был добавлен компонент социального заражения после фиксации основной двухкомпонентной модели на этапе ранней проверки.

Сценарная роль модели отклика отделена от роли прогноза. Модель не претендует на воспроизведение фактического поведения аудитории конкретной конференции; она описывает фиксированную параметризуемую структуру отклика, относительно которой выполняется численное сравнение политик и вариантов программы. Чувствительность результатов сравнения к допущениям модели исследуется в эксперименте отдельно через варьирование параметрических осей, описывающих веса компонентов, размер аудитории и источник оценки популярности докладов. Эта проверка устойчивости --- самостоятельная задача, выполняемая в главе~\ref{ch:results}.


\section{Реализация управляющих воздействий: политики рекомендаций и оператор локальных модификаций программы}
\label{sec:3-4}

У организатора программы конференции имеется два качественно различных канала влияния на ожидаемое распределение посещений. Первый канал --- алгоритмический: формирование рекомендаций отдельным участникам в каждом временном слоте через выбранную политику. Второй канал --- структурный: локальная модификация привязки докладов к временным слотам с сохранением общего состава программы. В программном комплексе оба канала реализованы как взаимозаменяемые компоненты, подключаемые к симулятору через единый контракт; настоящий подраздел описывает реализацию каждого из них и обоснование выбора конкретных дизайнов.

\paragraph{Семейство политик рекомендаций.}
В сравнении участвует семейство из четырёх политик, отвечающее минимальной достаточности для содержательного сравнительного анализа.

Контрольная политика не предъявляет участнику никакой рекомендации. В функции полезности участника компонент следования рекомендации тождественно равен нулю на всех докладах слота, и распределение выбора определяется только содержательной близостью и компонентом социального заражения. Эта политика выполняет роль точки отсчёта: при любых значениях параметрических осей её результаты задают опорную картину распределения посещений, формирующегося без алгоритмического вмешательства. Без такой точки отсчёта сравнение остальных политик теряло бы интерпретируемость, поскольку любой эффект политики был бы неотличим от эффекта самой по себе аддитивной функции полезности.

Политика по релевантности формирует упорядоченный список рекомендаций по убыванию содержательной близости описания доклада к профилю участника. Сигналы загрузки залов и популяции она не использует. Эта политика моделирует наиболее распространённый в практических рекомендательных системах малого масштаба класс подходов, в том числе тот, который реализован в существующем сервисе рекомендаций партнёра. Она служит ориентиром для оценки эффекта учёта системных ограничений в более сложных политиках.

Политика с учётом загрузки залов сочетает ранжирование по релевантности с явным учётом текущей заполненности залов. К скорингу кандидатов добавляется компонент, монотонно убывающий по доле занятых мест в зале, в котором проходит доклад; залы с заполненностью выше заданного жёсткого порога исключаются из выдачи. При недостаточном числе допустимых кандидатов жёсткое отсечение снимается, чтобы выдача сохраняла ненулевую длину. Эта политика реализует простое содержательное правило, выдвигаемое организатором мероприятия: при прочих равных направлять участников в менее загруженные залы и не направлять в почти переполненные. В ней сосредоточен основной канал учёта вместимости, вынесенный из функции полезности участника по соображениям, изложенным в подразделе~\ref{sec:3-3}.

Политика на основе языковой модели формирует выдачу через вызов внешней языковой модели, которой передаётся текстовый профиль участника и текстовые описания кандидатов слота с просьбой вернуть упорядоченный список идентификаторов. Эта политика реализует четвёртый качественно различный класс --- семантическое ранжирование без явно заданной функции близости. От политики по релевантности она отличается тем, что строит ранжирование через семантический разбор текстов и не опирается на фиксированную косинусную близость векторных представлений; от политики с учётом загрузки залов --- отсутствием явного канала ограничений вместимости. Включение её в семейство проверяет, даёт ли семантическое ранжирование заметный выигрыш относительно простой косинусной близости при равенстве остальных условий. В реализации использован вариант без передачи в запрос состояния загрузки залов: это сохраняет вынос ограничения вместимости на уровень соответствующей политики семейства и обеспечивает сопоставимость с политикой по релевантности по составу учитываемых сигналов. Результаты вызовов кэшируются по идентификаторам участника и слота для воспроизводимости и контроля повторных запусков.

Минимальная достаточность четырёх политик обоснована тем, что каждая представляет качественно различный класс подхода: отсутствие алгоритмического вмешательства, классическое ранжирование по релевантности, учёт системных ограничений вместимости, семантическое ранжирование на основе языковой модели. Обучаемые политики на основе графовых нейронных сетей или обучения с подкреплением исключены из семейства из-за отсутствия обучающего материала, что обсуждалось в главах~\ref{ch:domain} и~\ref{ch:method}. Алгоритмы переранжирования с управлением разнообразием выдачи рассматривались отдельно: задача управления разнообразием не сводится к учёту вместимости и не входит в центральный вопрос работы --- сравнение политик по показателям риска перегрузки.

Все политики реализованы как взаимозаменяемые компоненты с единым контрактом подключения к симулятору. Это даёт прозрачное сравнение политик в одной экспериментальной обстановке и возможность расширения семейства без изменения симулятора и компонентов оценки. Политика на основе языковой модели подключается через тот же контракт, что и три остальные, с единственным отличием --- стоимостью и латентностью вызовов; этот аспект учитывается в дизайне эксперимента, не в логике симулятора.

\paragraph{Оператор локальных модификаций программы.}
Оператор локальной модификации программы реализует структурный канал управления. Он принимает на вход исходную программу и возвращает конечное множество допустимых вариантов программы, полученных перестановками пар докладов между разными временными слотами одного дня. Модификация состоит в том, что у двух выбранных докладов меняется временной слот; привязка к залу при этом сохраняется, состав докладов программы остаётся неизменным. Полученные варианты подаются в симулятор так же, как исходная программа, и образуют отдельную ось параметрической конфигурации.

При выборе формы оператора рассматривался ряд альтернатив. Полная задача оптимизации программы средствами целочисленного программирования, разработанная в литературе по conference scheduling --- в работах Вангервена и соавторов, Резайнии и соавторов, Пилявского и соавторов, --- была отвергнута по двум причинам. Во-первых, такие постановки опираются на допущение известных или восстановленных опросом преференций участников; в условиях настоящей работы соответствующий канал данных недоступен. Во-вторых, и это главное, задача работы --- не нахождение оптимального расписания, а сценарный анализ устойчивости программы к её локальным изменениям. Применение полной оптимизации в этой роли было бы методологически некорректным: сравнение исходной программы с глобально оптимизированной не отражает интересующего вопроса об эффекте небольших, реалистичных для организатора перестановок.

Альтернативные локальные операторы --- обмен блоков по нескольким докладам, перестановки между залами внутри одного слота, эвристический локальный поиск по риску перегрузки --- также рассматривались. Обмен блоков увеличивает мощность изменения и приближает оператор к полной оптимизации, что снижает интерпретируемость результата как <<небольшой реалистичной перестановки>>. Перестановки между залами внутри одного слота не меняют распределение посещений по слотам и интересны как отдельный самостоятельный канал, выходящий за рамки настоящей работы. Эвристический локальный поиск по риску перегрузки превращает оператор в оптимизатор и противоречит роли структурной оси сценарного эксперимента.

Принятая форма оператора --- попарная перестановка двух докладов между разными временными слотами одного дня. Это минимальная локальная модификация, не выходящая за пределы одного дня программы, понятная организатору как реалистичный приём управления программой и дающая равномерную сопоставимость вариантов. Бюджет оператора задаётся числом возвращаемых модификаций; при превышении бюджета модификации выбираются случайной подвыборкой с фиксированным случайным зерном на уровне точки гиперкуба, что даёт воспроизводимость и одинаковость множества вариантов для всех политик внутри одной точки.

Ограничение на конфликты спикеров реализовано как жёсткая проверка: модификация, в результате которой один спикер оказывается одновременно в двух разных временных слотах одного дня, исключается из множества возвращаемых вариантов. Проверка применяется ко всем кандидатам и предшествует случайной подвыборке. На программах, в исходных данных которых поле спикеров отсутствует, проверка становится тривиальной: результаты в этом случае трактуются как отсутствие данных о возможных конфликтах, а не как отсутствие конфликтов в качестве свойства программы. Это уточнение существенно для интерпретации результатов на наборах данных без поля спикеров и фиксируется как одно из условий применения оператора.

Исходная программа включается в множество вариантов как отдельная контрольная точка с фиксированным индексом. При сравнении эффекта структурных модификаций исходная программа всегда присутствует на той же оси параметрической конфигурации, что и её модификации, и попарное сравнение модификаций с исходной программой выполняется без отдельной служебной логики.

Алгоритмический и структурный каналы управления реализованы как ортогональные оси параметрической конфигурации: при каждом запуске эксперимента симулятор получает сначала вариант программы, выбранный по соответствующей оси, затем политику рекомендаций, выбранную по другой оси, и обрабатывает их совместно. Это позволяет в главе~\ref{ch:results} раздельно оценить эффект алгоритмической логики и эффект структуры программы на показатели риска и устойчивости.


\section{Реализация LLM-агентского симулятора}
\label{sec:3-5}

В составе комплекса реализован второй симулятор отклика участника на основе большой языковой модели. Он не заменяет параметрический симулятор и не служит самостоятельным результатом работы; его роль --- взаимная проверка выводов параметрического симулятора о сравнительной устойчивости политик на репрезентативном подмножестве конфигураций. Параметрический и агентский симуляторы образуют не два независимых в строгом смысле источника отклика, а два источника отклика с принципиально разной формой модели выбора: в параметрическом --- фиксированная функция полезности с явными весами компонентов, в агентском --- свободная форма решения через цепь рассуждений языковой модели по тексту профиля и контекста. Совпадение или расхождение выводов двух источников при общем входе и общем контракте политики служит диагностическим сигналом о чувствительности результата к форме модели выбора. Настоящий подраздел описывает реализацию агента, состав передаваемого ему контекста и обоснование принятых ограничений.

\paragraph{Состав агента.}
Агент реализован в минимальной архитектуре <<профиль --- память --- действие>>. Профиль агента --- это текстовое описание участника, согласованное с тематикой конференции; в нём фиксируются интересы, уровень квалификации, предпочтения по форматам выступлений и стилевые ожидания. Память агента --- это упорядоченный список докладов, посещённых им в текущем эпизоде до момента принятия очередного решения; при формировании очередного решения список передаётся агенту в текстовой форме с указанием идентификаторов и заголовков посещённых докладов. Действие агента --- выбор одного из доступных в текущем временном слоте докладов или явный отказ от посещения слота; решение возвращается в структурированной форме с кратким текстовым обоснованием.

В текущем эпизоде агенты обрабатываются последовательно по слотам и последовательно внутри одного слота. Каждый следующий агент в слоте наблюдает обновлённое состояние, отражающее решения тех агентов, что были обработаны раньше; так формируется кумулятивное распределение посещений, согласованное с параметрическим симулятором. После завершения слота агент фиксирует посещение в собственной памяти, и эта запись становится доступной ему при принятии решений в последующих слотах того же эпизода.

Эта архитектура --- стандартная минимальная сборка для агентских симуляторов пользователей рекомендательных систем, использованная в работе Парка и соавторов о генеративных агентах и в средах Agent4Rec, OASIS и AgentSociety, обзор которых приведён в главе~\ref{ch:domain}. Расширения минимальной архитектуры --- пятифакторная модель личности, граф знакомств между агентами, рефлексивная переработка опыта, межслотовое общение участников через общий пул впечатлений --- рассматривались на этапе проектирования и оставлены за пределами объёма работы. Каждое такое расширение добавляет поведенческие элементы, которые сами требуют отдельной проверки на согласованность; роль LLM-симулятора в настоящей работе не в более точном моделировании аудитории, а в независимом источнике отклика для перекрёстной проверки выводов параметрического симулятора на общем подмножестве конфигураций.

\paragraph{Состав передаваемого контекста и паритет с параметрическим симулятором.}
Существенная часть проектирования LLM-агента --- определение того, что именно передаётся в запрос модели в качестве контекста принятия решения. Принципом, по которому делалось это определение, является паритет каналов с параметрическим симулятором: по каждому каналу влияния на выбор участника решается, передавать соответствующую информацию агенту или нет, и решение это принимается симметрично с реализацией соответствующего канала в функции полезности параметрического симулятора.

Текстовый профиль и история сегодняшних посещений передаются всегда: они соответствуют каналу содержательной релевантности и динамике индивидуального предпочтения по ходу эпизода. Упорядоченный список рекомендаций политики передаётся всегда: он соответствует каналу следования рекомендации и позволяет сравнивать политики в LLM-симуляторе на тех же основаниях, что и в параметрическом. Состояние загрузки залов агенту не передаётся: по принятому в подразделе~\ref{sec:3-3} разделению каналов ограничение вместимости трактуется как ответственность политики, и передача его в промпт агента нарушила бы это разделение, дав агенту прямой доступ к информации, которая в параметрическом симуляторе ему недоступна.

Социальный сигнал --- число участников, уже выбравших каждый из докладов слота к моменту принятия решения текущим агентом --- передаётся в промпт условно. При выключенном канале социального заражения, когда соответствующий вес в функции полезности параметрического симулятора равен нулю, социальный сигнал не передаётся и в LLM-симулятор. При ненулевом весе сигнал передаётся в виде кратких числовых сводок, и в системный промпт добавляется указание учитывать его как умеренный или заметный социальный фактор в зависимости от значения веса. Это решение выровнено с реализацией социального заражения в параметрическом симуляторе: оба источника отклика работают с одной и той же семантикой социального сигнала. Конкретные числовые значения сигнала формируются эндогенно внутри каждого симулятора по ходу эпизода и могут различаться между симуляторами; различие трактуется как следствие различной формы модели выбора, а не как несоответствие реализаций.

\paragraph{Бюджет и ограничения масштаба.}
LLM-симулятор требует существенно больших вычислительных ресурсов и материальных затрат, чем параметрический. Каждое решение каждого агента требует отдельного вызова внешней языковой модели; стоимость и время выполнения растут пропорционально произведению размера аудитории, числа слотов и числа сравниваемых политик. Применить LLM-симулятор на полном плане параметрического эксперимента нецелесообразно ни по времени, ни по бюджету.

В принятой схеме LLM-симулятор работает на репрезентативном подмножестве конфигураций, отобранных по принципу равномерного покрытия параметрического пространства, и на сокращённой синтетической аудитории. Внутри одной точки этого подмножества все четыре политики прогоняются на одной и той же синтетической аудитории --- реализация принципа общих случайных чисел, аналогичная параметрическому симулятору. Контроль стоимости реализован как явное ограничение сверху на суммарные затраты прогона: при достижении заданного предела прогон останавливается, частичные результаты сохраняются, и факт превышения предела фиксируется в метаданных эксперимента. Все вызовы языковой модели выполняются асинхронно с ограничением числа одновременных запросов, для каждого вызова фиксируется фактическая стоимость и время выполнения.

\paragraph{Принципиальная роль второго симулятора.}
Назначение LLM-симулятора в работе --- не получить более достоверную модель аудитории, а проверить выводы параметрического симулятора альтернативным источником отклика. Это совпадает с критической линией литературы по агентским симуляторам, обсуждённой в главе~\ref{ch:domain}: в работе Лароожа и Тёрнберга показано, что индивидуальная точность LLM-агентов невелика, а правдоподобие их поведения не эквивалентно операциональной валидности модели. Поэтому LLM-симулятор не служит источником самостоятельных абсолютных утверждений о поведении аудитории. Все его выводы рассматриваются только в сопоставлении с выводами параметрического симулятора на общем подмножестве конфигураций; совпадение ранжирований политик между двумя симуляторами трактуется как взаимная проверка двух источников отклика с разными формами модели выбора.


\section{Реализация экспериментального дизайна и подготовки данных для валидации}
\label{sec:3-6}

Реализация плана эксперимента отделена от логики симуляторов и оформлена в виде самостоятельного модуля, формирующего набор параметрических конфигураций по заданным осям и обеспечивающего разделение источников случайности между ними. Настоящий подраздел описывает реализацию плана эксперимента и обоснование принятых решений.

\paragraph{Латинский гиперкуб как метод покрытия пространства параметров.}
Пространство параметрических конфигураций задано шестью осями: множитель вместимости залов, источник оценки популярности докладов, вес компонента следования рекомендации, вес компонента социального заражения, размер синтетической аудитории, индекс варианта программы. Третий компонент функции полезности --- содержательная релевантность --- получает дополняющий вес из симплексного ограничения и собственной оси не образует. Параметр стохастичности softmax (температура) и базовая вероятность отказа от посещения слота в плане эксперимента не варьируются: они рассматриваются как фиксированные настройки симулятора, а не как оси основного сценарного эксперимента.

Обоснование выбора латинского гиперкуба как плана эксперимента приведено в главе~\ref{ch:method}. В реализации план содержит пятьдесят точек по шести осям; внутри каждой точки перебираются все четыре политики семейства, и каждая комбинация <<точка $\times$ политика>> повторяется на трёх независимых случайных зёрнах. Этот объём согласован с временным бюджетом полного параметрического прогона и даёт достаточную статистическую базу для попарного сравнения политик внутри точки в режиме общих случайных чисел. Полный факторный перебор по тому же пространству при разумной дискретизации даёт количество комбинаций, превышающее временной бюджет в десятки раз; глобальные методы оценки чувствительности избыточны для задачи попарного сравнения политик внутри точки конфигурации.

\paragraph{Симплексное ограничение весов и сбалансированность дискретных осей.}
Веса трёх компонентов функции полезности участника связаны симплексным ограничением: их сумма равна единице. При генерации точек гиперкуба не все случайные комбинации значений соответствующих координат удовлетворяют этому ограничению, и точки, нарушающие его, исключаются. Реализован циклический алгоритм генерации блоков точек с последующей фильтрацией: каждый блок латинского гиперкуба порождается детерминированно от собственного под-зерна, образованного от главного зерна эксперимента; точки, нарушающие симплекс, отбрасываются; цикл продолжается до набора требуемого числа допустимых точек. Такой порядок воспроизводим и не зависит от порядка генерации блоков.

После накопления требуемого числа точек выполняется проверка сбалансированности дискретных осей --- индекса варианта программы, размера аудитории и источника оценки популярности докладов. Поскольку рандомизация в латинском гиперкубе сама по себе не гарантирует сбалансированности дискретных осей при небольшом числе точек, при недостаточном представлении какого-либо уровня дискретной оси выполняется замена случайной избыточно представленной точки на принудительно представленную точку соответствующего уровня. Замены фиксируются в метаданных плана и не нарушают свойство равномерного покрытия гиперкуба, поскольку выполняются на оси, дискретной по построению.

\paragraph{Подмножество для агентского симулятора.}
Агентский симулятор работает на репрезентативном подмножестве из двенадцати точек, отобранных из полных пятидесяти. Принципом отбора служит равномерное покрытие пространства: подмножество должно охватывать как крайние, так и центральные значения каждой из осей конфигурации. Этому требованию отвечает алгоритм максимизации минимального попарного расстояния --- последовательный отбор точек, при котором каждая следующая точка максимально удалена от уже отобранных в нормированном пространстве осей. Альтернативные методы --- отбор по кластерам, ручной отбор по краям и центру, равномерное проектирование на каждую ось --- рассматривались на этапе проектирования; алгоритм максимизации минимального расстояния выбран как метод, формализуемо реализующий требование равномерного покрытия без участия эксперта.

Дополнительно к алгоритму отбора применяется принудительное включение точки с исходным вариантом программы: это сохраняет в подмножестве контрольную точку структурного канала управления и поддерживает интерпретацию исходной программы как опорной для сравнения её локальных модификаций.

\paragraph{Общие случайные числа.}
Принцип общих случайных чисел применяется на уровне точки гиперкуба. Внутри одной точки случайное зерно отбора синтетической аудитории и случайное зерно оператора локальных модификаций программы фиксированы; они одинаковы для всех политик и для всех реплик внутри этой точки. Все политики и все реплики работают с одной и той же синтетической аудиторией и с одной и той же реализацией варианта программы. Допустимый источник случайности при репликации --- собственное зерно симулятора --- изолирован от первых двух и варьируется между репликами. Внутри самого симулятора разделены поток случайности окончательного выбора участника и поток случайности стохастической части политики; это делает граничное поведение модели строго воспроизводимым.

Такое разделение существенно для содержательной интерпретации сравнения политик. Без общих случайных чисел шум, вносимый разной выборкой аудитории при сравнении двух политик, мог бы маскировать или искусственно усиливать различие между ними. С общими случайными числами это различие отражает только различие в логике политик при одинаковом внешнем условии; именно это и нужно для попарного сравнения внутри точки конфигурации.

\paragraph{Ограничение запуска политики на основе языковой модели.}
Стоимость запуска политики на основе языковой модели существенно превосходит стоимость запуска остальных политик семейства. Полный охват ею всех пятидесяти точек параметрического плана эксперимента не укладывается в практический временной и денежный бюджет работы. В принятой схеме политика на основе языковой модели запускается в параметрическом симуляторе только на тех же двенадцати точках, что и подмножество для агентского симулятора. Это даёт перекрёстную сопоставимость двух симуляторов на общих двенадцати точках и сохраняет возможность сравнения трёх остальных политик семейства на полных пятидесяти точках. Решение фиксируется в метаданных эксперимента и в разделе ограничений работы.

\paragraph{Формат результатов.}
Результаты единичных итераций сохраняются в табличном формате с явным указанием идентификатора точки гиперкуба, значений всех параметрических осей конфигурации, идентификатора политики, номера реплики и значений всех показателей групп, описанных в подразделе~2.2.3. Формат напрямую совместим с модулем постобработки, который вычисляет попарные сравнения политик, оценки устойчивости при варьировании одной из осей и эффект структурных модификаций программы. Постобработка не зависит от внутренней логики симуляторов: она работает как функция, отображающая таблицу результатов в таблицы попарных сравнений и характеристик чувствительности. Численные результаты этой обработки и их интерпретация представлены в главе~\ref{ch:results}.


\section{Использованные данные и верификация работоспособности}
\label{sec:3-7}

Программный комплекс работает на наборе данных, состоящем из программы основной конференции, контролируемых синтетических программ для проверки граничных свойств реализации и пред-сгенерированной популяции синтетических профилей участников. Настоящий подраздел описывает назначение каждого типа данных и реализованные процедуры проверки работоспособности комплекса.

\paragraph{Программа основной конференции.}
В качестве основной конференции для численного эксперимента использована программа Mobius 2025 Autumn --- мероприятия сообщества JUG, посвящённого мобильной разработке. Программа охватывает два дня и содержит сорок докладов, распределённых по шестнадцати временным слотам: в двенадцати параллельных слотах используются три зала, в четырёх пленарных --- один. Номинальная вместимость каждого зала составляет сто посадочных мест. Выбор этой программы обусловлен сочетанием трёх факторов: наличие у партнёра реальной программы с фактическими данными о залах и докладах, размер программы, достаточный для содержательного сценарного анализа без чрезмерной вычислительной стоимости, и согласованность тематики программы с тематикой подготовленной популяции синтетических профилей участников.

Метаданные программы --- текстовые описания докладов, тематические признаки и состав спикеров --- используются на разных этапах работы комплекса. Текстовые описания служат основой для формирования векторных представлений докладов, используемых в политиках рекомендаций и в функции полезности модели поведения участника. Тематические признаки и сопутствующие структурные характеристики программы используются для вычисления оценки популярности докладов в одном из вариантов источника популярности. Состав спикеров используется оператором локальных модификаций программы для проверки ограничения на конфликты спикеров.

\paragraph{Синтетическая популяция участников.}
Аудитория задаётся синтетическим порождающим процессом. Для основной конференции подготовлена популяция из ста индивидуальных профилей участников. Каждый профиль --- это текстовое описание интересов, уровня квалификации, тематических предпочтений и стилевых ожиданий, согласованное с тематикой основной конференции, и нормированное векторное представление этого описания, используемое в функции полезности модели поведения. Синтетические аудитории конкретных размеров формируются отбором подмножеств из этой популяции с фиксированным случайным зерном, образуемым по описанному в подразделе~\ref{sec:3-6} правилу.

Пред-сгенерированная популяция с последующим отбором подмножеств выбрана в противовес генерации новой популяции на каждом запуске. Это сокращает стоимость подготовки данных, даёт воспроизводимость и согласованность аудиторий между прогонами, поддерживает корректное применение принципа общих случайных чисел на уровне точки гиперкуба. Разнородность популяции --- профили разной квалификации и разных тематических подобластей --- фиксируется на этапе подготовки и не зависит от случайного зерна отбора подмножества.

\paragraph{Контролируемые синтетические программы.}
Помимо программы основной конференции в составе комплекса подготовлены три контролируемые синтетические программы --- toy-микроконференции, --- назначение которых отлично от роли основной конференции. Эти программы используются для верификации граничных свойств реализации и не участвуют в основном сценарном эксперименте.

Первая toy-микроконференция содержит один временной слот с двумя параллельными залами и двумя докладами и предназначена для проверки базовых свойств функции полезности и распределения выбора при контролируемых параметрах: при нулевом весе компонента следования рекомендации все политики дают близкие распределения посещений, при ненулевом весе различие политик выражено и монотонно растёт по этому весу. Вторая toy-микроконференция содержит два временных слота с такой же структурой каждого слота и предназначена для проверки кумулятивного эффекта между слотами в LLM-агентском симуляторе: на одной слотовой программе агент с памятью сегодняшних посещений неотличим от агента без памяти, и для содержательной проверки этого канала требуется как минимум двухслотовая программа. Третья контролируемая программа содержит специально внесённый конфликт спикеров --- одного спикера в двух разных временных слотах одного дня --- и предназначена для проверки того, что оператор локальных модификаций программы корректно отвергает варианты, нарушающие соответствующее ограничение.

\paragraph{Граничная верификация модели.}
Граничная верификация модели реализована как блокирующий фильтр перед содержательными выводами в соответствии с четырьмя свойствами, сформулированными в главе~\ref{ch:method}: исчезающий риск перегрузки при избыточной номинальной вместимости, монотонный рост риска при её уменьшении, неразличимость политик при нулевом весе компонента следования рекомендации и максимальное различие политик при единичном весе того же компонента.

Проверка реализована как набор автоматизированных тестов, каждый из которых выполняется на минимально достаточном размере конфигурации с усреднением по нескольким случайным зёрнам. Тесты предусмотрены для запуска при проверке изменений ядра комплекса; при провале хотя бы одного из них переход к содержательным результатам блокируется до устранения причины. Такая процедура соответствует канону верификации имитационных моделей [Sargent, 2013; Robinson, 2014], упомянутому в главе~\ref{ch:domain}. На текущей редакции комплекса все четыре граничных свойства выполняются --- соответствующая часть автоматизированных тестов проходит без отказов; численные характеристики граничной верификации представлены в главе~\ref{ch:results}.

\paragraph{Юнит-тестовое покрытие модулей.}
Помимо граничной верификации в составе комплекса реализованы юнит-тесты на отдельные модули: вычисление компонентов функции полезности, вычисление показателей загрузки и риска, корректность работы политик, корректность работы оператора локальных модификаций программы и проверки ограничения на конфликты спикеров, корректность работы генератора плана эксперимента, корректность правил образования случайных зёрен. Юнит-тесты не пересекаются с граничной верификацией: первые проверяют корректность отдельных компонентов в изоляции, вторая --- корректность модели в целом на уровне эмерджентных свойств. Два уровня тестирования вместе позволяют проверить, что отдельные модули комплекса работают корректно, а собранная из них модель ведёт себя ожидаемо в граничных режимах.


\section{Технологический стек и использование систем искусственного интеллекта}
\label{sec:3-8}

\paragraph{Технологический стек.}
Технологическая сложность реализованного комплекса определяется не выбором конкретных библиотек, а инженерными решениями, описанными в подразделах~\ref{sec:3-3},~\ref{sec:3-5} и~\ref{sec:3-6}: разделением каналов влияния политики и состояния системы; двусимуляторной архитектурой моделирования отклика с паритетом каналов между двумя симуляторами; реализацией принципа общих случайных чисел через явное разделение источников случайности на трёх уровнях (отбор аудитории, оператор локальных модификаций программы, симулятор) и двух потоков случайности внутри симулятора (выбор участника, стохастическая часть политики); оператором локальных модификаций программы с жёсткой проверкой ограничений; асинхронной обработкой вызовов внешних языковых моделей с контролем стоимости и кэшированием. Конкретный технологический стек собран из типовых библиотек экосистемы Python; нестандартная сложность вынесена в перечисленные архитектурные решения.

Программный комплекс реализован на языке Python. Численные вычисления выполняются средствами библиотек NumPy и SciPy; SciPy используется в том числе для генерации точек латинского гиперкуба. Обучение вспомогательной модели предпочтений, использованной в одном из вариантов оценки релевантности, выполнено средствами библиотеки scikit-learn методом градиентного бустинга на решающих деревьях. Векторные представления описаний докладов и текстовых профилей участников получены однократным запросом модели многоязычных эмбеддингов \texttt{intfloat/multilingual-e5-small}; результаты сохранены в бинарных файлах с привязкой к идентификаторам.

Доступ к большим языковым моделям, использованным в политике на основе языковой модели и в LLM-агентском симуляторе, реализован через единый совместимый с OpenAI протокол с маршрутизацией через сервис OpenRouter. В политике на основе языковой модели использована модель \texttt{openai/gpt-4o-mini}; в LLM-агентском симуляторе для основного прогона на репрезентативном подмножестве конфигураций использована экономичная модель того же семейства; конкретное имя модели зафиксировано в метаданных каждого запуска. Все вызовы выполняются асинхронно с ограничением числа одновременных запросов и обработкой ошибок внешнего сервиса; для каждого вызова фиксируется фактическая стоимость и время выполнения. Реализован дисковый кэш ответов политики на основе языковой модели --- повторный запуск той же конфигурации не оплачивается дважды; ответы LLM-агентов в текущей реализации не кэшируются.

Управление виртуальным окружением и зависимостями выполнено с использованием утилиты uv. Хранение версий кода и результатов экспериментов организовано через систему контроля версий git с ведением единой основной ветки. Метаданные каждого запуска эксперимента --- точные значения параметров, используемые случайные зёрна, время выполнения, фактическая стоимость в случае использования внешних языковых моделей --- сохраняются в составе результатов; этого достаточно для воспроизведения параметрической части эксперимента и трассировки условий LLM-запусков.

\paragraph{Использование систем искусственного интеллекта при подготовке работы.}
В соответствии с требованиями к выпускной квалификационной работе, действующими в Университете ИТМО, в составе работы фиксируется применение систем искусственного интеллекта при её подготовке.

Системы искусственного интеллекта применялись в работе как инструмент моделирования аудитории внутри программного комплекса, как инструмент разработки и как инструмент подготовки текстов и анализа источников. Первое применение --- внутренний компонент разрабатываемой системы и описано в подразделах~\ref{sec:3-4} и~\ref{sec:3-5}: большая языковая модель работает как ранжирующий компонент в одной из политик семейства и как источник отклика в агентском симуляторе. Это применение --- часть предмета работы и подвергается тем же процедурам верификации, что и остальные модули комплекса.

Применение систем искусственного интеллекта в качестве инструмента разработки и инструмента подготовки текстов выполнялось вспомогательно. Языковые модели использовались для генерации фрагментов кода по техническим заданиям, формулируемым автором, для предложения вариантов оформления отдельных текстовых разделов, для пояснения деталей внешних библиотек и для поиска и реферирования научных источников при подготовке обзора предметной области в главе~\ref{ch:domain}. Полученные фрагменты во всех случаях проходили содержательный контроль автора: для кода --- функциональное тестирование и интеграцию в комплекс с прохождением реализованных юнит-тестов и граничной верификации; для текстов --- редакторскую проверку на соответствие постановке работы и на стилистическое единство с остальной частью текста; для источников --- независимую проверку существования и содержания каждой цитируемой работы по первичным источникам. Принятие архитектурных решений, постановка задачи, выбор методов, интерпретация полученных результатов и итоговая ответственность за работу принадлежат автору.


\section{Выводы по главе}
\label{sec:3-9}

В главе представлена реализация программного комплекса для сценарного анализа программы конференции в условиях неопределённости спроса и отсутствия данных о фактической посещаемости. Состав комплекса отвечает постановке главы~\ref{ch:method} и собран из взаимозаменяемых компонентов: внутреннее представление программы конференции, синтетическая аудитория, модель поведения участника, семейство политик рекомендаций, оператор локальных модификаций программы, два варианта симулятора отклика, генератор плана эксперимента, модуль показателей и средства формирования сравнительного отчёта.

Структуру комплекса определяют три инженерных решения, выполненных в ходе работы. Первое --- разделение каналов влияния политики и состояния системы на функцию полезности участника, при котором ограничение вместимости вынесено из модели поведения на уровень одной из политик семейства; благодаря этому требования граничной верификации выполняются по построению. Второе --- двусимуляторная архитектура моделирования отклика: параметрический симулятор работает на полном плане эксперимента, а агентский симулятор на основе языковой модели --- на репрезентативном подмножестве конфигураций для перекрёстной проверки выводов. Третье --- план эксперимента на основе латинского гиперкуба с общими случайными числами внутри точки конфигурации; это даёт равномерное покрытие пространства параметров и низкошумное попарное сравнение политик.

Сравнительные результаты численного эксперимента на программе Mobius 2025 Autumn, попарные оценки устойчивости политик при варьировании параметрических осей, эффект локальных модификаций программы относительно исходной и согласованность ранжирования политик между параметрическим и LLM-агентским симуляторами на общем подмножестве конфигураций представлены в главе~\ref{ch:results}.
